
\section{Statistical Analysis}

This section analyses the recorded badminton match
results using statistical methods. The objective is to
understand the team's performance across Singles,
Doubles, and Mixed Doubles formats.

The analysis considers the points scored by our team,
the points scored by the opponents, and the difference
between the two. Descriptive statistics are used to
identify general performance patterns and variations
in the recorded results.

% ==================================================
\subsection{Performance Parameters}
% ==================================================

The following parameters are used throughout the
analysis.

\begin{itemize}

    \item \textbf{Wins:} Points scored by our team.

    \item \textbf{Losses:} Points scored by the
    opposing team.

    \item \textbf{Net Result:} Difference between
    points scored by our team and the opponent.

\end{itemize}

For a match record, the net result is calculated as

\[
R_i = W_i - L_i
\]

where:

\begin{itemize}
    \item $W_i$ represents the points scored by our
    team in the $i$-th match.
    \item $L_i$ represents the points scored by the
    opponent in the $i$-th match.
    \item $R_i$ represents the net result.
\end{itemize}

A positive net result indicates that our team scored
more points, while a negative value indicates that
the opponent scored more points.

% ==================================================
\subsection{Descriptive Statistics}
% ==================================================

Descriptive statistics provide a summary of the
recorded performance in each match format.

The following statistical measures are considered.

\subsubsection{Mean}

The arithmetic mean represents the average number of
points scored per recorded match.

For the points scored by our team,

\[
\overline{W} =
\frac{1}{n}\sum_{i=1}^{n} W_i
\]

Similarly, the average points scored by the opponent
are given by

\[
\overline{L} =
\frac{1}{n}\sum_{i=1}^{n} L_i
\]

The average net result is

\[
\overline{R} =
\frac{1}{n}\sum_{i=1}^{n} R_i
\]

where $n$ is the number of valid records for the
respective match format.

\subsubsection{Standard Deviation}

The standard deviation measures the variation of
individual match results around their mean.

For our team's points,

\[
\sigma_W =
\sqrt{
\frac{1}{n}
\sum_{i=1}^{n}
(W_i-\overline{W})^2
}
\]

A lower standard deviation indicates that the
recorded scores are relatively close to the average.
A higher value indicates greater variation between
matches.

The same calculation can be applied to the opponent's
scores and the net results.

\subsubsection{Total Points}

The total points scored by both teams are calculated
as

\[
W_{\text{total}} =
\sum_{i=1}^{n} W_i
\]

\[
L_{\text{total}} =
\sum_{i=1}^{n} L_i
\]

The cumulative net result is

\[
R_{\text{total}} =
W_{\text{total}} - L_{\text{total}}
\]

These values provide an overall comparison of
scoring performance within each match format.

% ==================================================
\subsection{Comparison Across Match Formats}
% ==================================================

The three match formats are analysed independently
to understand how performance differs between
Singles, Doubles, and Mixed Doubles.

The comparison considers:

\begin{enumerate}
    \item Number of recorded matches.
    \item Total points scored by our team.
    \item Total points scored by opponents.
    \item Average points scored per match.
    \item Average net result.
    \item Standard deviation of net results.
\end{enumerate}

The comparison can be summarised using the following
table.

\begin{table}[H]
\centering
\caption{Summary of Statistical Measures}
\begin{tabular}{lccc}
\toprule
\textbf{Parameter} &
\textbf{Singles} &
\textbf{Doubles} &
\textbf{Mixed} \\
\midrule
Number of Matches & -- & -- & -- \\
Total Wins        & -- & -- & -- \\
Total Losses      & -- & -- & -- \\
Average Wins      & -- & -- & -- \\
Average Losses    & -- & -- & -- \\
Average Net Result& -- & -- & -- \\
Std. Deviation    & -- & -- & -- \\
\bottomrule
\end{tabular}
\end{table}

The table allows the three formats to be compared
using consistent statistical measures.

% ==================================================
\subsection{Date-wise Performance Analysis}
% ==================================================

Date-wise analysis is used to identify the variation
in performance over the observation period.

For each date, the net result is calculated using

\[
R_d = W_d - L_d
\]

where $W_d$ and $L_d$ represent the total points
recorded for our team and the opponent on date $d$,
respectively, for the selected match format.

The recorded dates can then be ranked according to
their net results.

A higher positive value indicates a stronger
scoring advantage, whereas a negative value indicates
a scoring disadvantage.

This ranking helps identify dates with comparatively
strong and weak performances.

% ==================================================
\subsection{Graphical Analysis}
% ==================================================

Graphs are used to make the statistical patterns
easier to interpret.

The following visualisations are considered for
each match format.

\subsubsection{Points Scored Over Time}

A time-series plot is used to represent the points
scored by our team and the opponent across the
recorded dates.

The horizontal axis represents the date, while the
vertical axis represents the number of points.

This plot helps identify changes in scoring
performance over the observation period.

\subsubsection{Net Result Over Time}

A separate plot represents the net result for each
recorded date.

A horizontal reference line at zero separates
positive and negative net results.

The average net result can also be represented as
a horizontal reference line to compare individual
dates against the overall average.

\subsubsection{Distribution of Net Results}

A distribution plot can be used to understand how
frequently different net-result values occur.

This helps determine whether the results are
concentrated around the average or spread across a
wide range of values.

% ==================================================
\subsection{Interpretation}
% ==================================================

The statistical measures and graphical
representations provide complementary views of
performance.

The mean describes the average scoring performance,
while the standard deviation indicates how much
individual results vary.

The cumulative net result gives an overall scoring
comparison, whereas date-wise rankings identify
individual periods of relatively strong or weak
performance.

Together, these measures establish a statistical
basis for comparing the three match formats and
provide the historical data required for the
next-game prediction model discussed in the
following section.

